% Modernised reading — fonds Alexandre Grothendieck, shelfmark 21, pages 1-28.
%
% This is an INTERPRETATION, not a transcription. It re-reads the pages in
% today's notation and vocabulary, states the hypotheses the manuscript leaves
% implicit, and drops the critical apparatus so that the mathematics reads
% straight through. Everything the brackets and underlines carried moves into a
% footnote. It is held to being mathematically correct as it stands; where the
% manuscript is loose or mistaken, the true statement is what appears here and
% the footnote says what the page has. In any dispute of substance the
% transcription governs: cite batch-01/02.fr.tex.
%
% Pass: Opus 5 (claude-opus-5), 2026-09-19 — first pass, unchecked against the
% pages by a human. The model was chosen explicitly for this pass.
%
% Scope: one document for the shelfmark, its two batches read as one argument.
% Twenty-four of the twenty-eight leaves carry mathematics; pages 1 and 14 are
% kraft wrappers, pages 2 and 12 blank versos.
%
% Spine. The folder asks ONE question and answers it three times, in three
% categories of increasing looseness: WHEN IS AN OBJECT ON X ALREADY KNOWN FROM
% ITS RESTRICTION TO X - Y, Y CLOSED? The answer is always the same shape — a
% vanishing condition on local cohomology along Y, checkable point by point —
% and the folder's work is to make that shape survive the loss of additivity.
%
%   pp. 3-13   THE TYPESCRIPT, « 2. Généralités sur certains foncteurs de
%              modules », Prop. 2.1 to Remarques 2.17. The additive toolkit:
%              (a) a cohomological functor vanishes iff it vanishes on a
%              generating family (2.1-2.5); (b) a left exact contravariant
%              functor on I-power-torsion modules is Hom(-, H) for a single
%              module H = colim T(A_n) (2.6-2.11); (c) consequences — the
%              vanishing criterion in terms of H (2.12, 2.13), Ass H = the
%              generic points (2.14), H injective iff T exact (2.15, 2.16),
%              and the sheaf-theoretic generalisation (2.17).
%   pp. 15-20  SHEAVES OF SETS, in his hand. (P1) and (P2): F -> i_* i^* F
%              injective, resp. bijective. Localisation Loc_x(X) at a point,
%              the pointwise form of the conditions, Prop. 3 (global iff
%              pointwise on Y), Cor. 4 (stability), Lemme 5 (sections over the
%              localisation are a filtered colimit), and the noetherian
%              induction.
%   pp. 21-28  FIBERED CATEGORIES, in his hand. Déf. 6: P1/P2/P3 = the
%              restriction functor is faithful / fully faithful / an
%              equivalence. Translation into the sheaves Hom(f,g); the induced
%              fibered category F_(x); Prop. 7 (global iff pointwise);
%              Cor. 8 (stability); the proofs; Cor. 10 (the conditions ARE the
%              vanishing of H^i_Y); and the closing special case of torsors,
%              Prop. 9, where H^0_Y and H^1_Y are the kernel and cokernel of
%              F -> i_* i^* F.
%
% What only a folder-wide reading can say. Three links, none of them stated on
% any one leaf, and the first is the folder's hinge:
%   - the module H^i = colim_m T^i(A_m) of 2.12 IS the local cohomology of the
%     manuscript half. The typescript builds it abstractly, as the object
%     representing a left exact functor, and never writes H^i_Y; the manuscript
%     writes H^i_Y and never says where it comes from. Read together, the
%     typescript is the construction of the very object the manuscript's
%     criterion is stated in.
%   - the noetherian induction is ONE move used twice in two categories: in 2.3
%     a maximal subobject admitting a composition series is shown to be the
%     whole object; on pp. 19 and 26 a maximal point of Y carries the induction
%     over the closed subsets of Y. The typescript proves the abstract form.
%   - the marginal N.B. of p. 15 — « regarder aussi la propriété : surjective »
%     — is answered eleven leaves later. Surjectivity at the level of OBJECTS,
%     not of sections, is exactly the third condition P3 of Déf. 6, and it is
%     what the whole of pp. 24-27 is spent on.
%
% NOTATION fixed for the folder. The two halves collide and are separated here:
%   U   — the typescript has no name for X - Y; pp. 15-20 write U, pp. 21-28
%         write a script U. WRITTEN HERE: U throughout.
%   F   — the typescript writes F for a general object of C (2.3, 2.4, 2.5);
%         pp. 15-28 write F for a sheaf and script F for the fibered category.
%         WRITTEN HERE: M, N for objects of C; F for a sheaf; script F for the
%         fibered category. Nothing is written F in the first half.
%   H   — the typescript's H (2.12 ff.) and the manuscript's H^i_Y are the same
%         object; the identification is made explicit in the spine section and
%         the two notations are kept apart in the body, because the folder
%         itself never joins them.
%   X_(x) — the set of generizations of x with the induced topology, his
%         Loc_x(X). X'_(x) = X_(x) - {x}.
% A is a commutative noetherian ring, I an ideal, A_n = A/I^{n+1}, X = Spec A,
% Y = V(I), U = X - Y, i : U -> X. Sheaf Hom is written script-H om.
%
% Substantive departures, each footnoted where it occurs:
%   - p. 3 (2.3): the leaf defines a system of generators by « il existe un
%     M in C », which is vacuous. It must be M in G, and is written so.
%   - p. 8: « modules de type fini sur C » where the argument needs A. Written
%     over A; the leaf's reading is footnoted, as the transcriber flagged it.
%   - p. 8: the same subcategory is written C^n and C_n on one leaf. C_n here.
%   - p. 9: Prop. 2.9 has a heading and no statement. The statement given here
%     is the one the surrounding paragraph forces, and is ours.
%   - pp. 9-10: Cor. 2.10 is cut by the leaf's end and 2.11 falls entirely in
%     the gap. 2.7 and 2.8 are attached to no legible statement and are
%     reconstructed only as far as the back-references allow.
%   - p. 13 (2.15): « En vertu de Krull » for what is the Artin-Rees lemma, and
%     the containment is written I^n N contains N cap I^m N where the argument
%     and the next display both need N cap I^m M.
%   - p. 27 (Cor. 10): the ranges are written i <= 1 (i <= 2, i <= 3) and are
%     one too high in each case; p. 28 states the correct ones for P1 and P2,
%     and the H^2_Y of the display on the same leaf confirms the third.
%   - p. 27: the illegible exponent on the right-hand H is 1, forced by
%     H^i_Y = R^{i-1} i_*; the reading is ours and is footnoted as such.
%   - p. 26: the proof glues an object over W with one over U along W cap U.
%     That is effectivity of descent for the covering {W, U} — the fibered
%     category must be a stack. Déf. 6 asks only for a fibered category.
%   - p. 22: the essential surjectivity clause is written with F(W) where the
%     argument needs F(V).
%   - p. 24: (ii) carries a capital X where the point x is meant.
%
% NUMBERING. His draft's, and disordered, and NOT reordered here: the
% typescript runs 2.1 to 2.17 with 2.7, 2.8, 2.9 and 2.11 carrying no legible
% statement; the manuscript runs Déf. 6, Prop. 7, Cor. 8 (p. 24), Cor. 10
% (p. 27), Prop. 9 (p. 28). The sections below follow the leaves, so Cor. 10 is
% read before Prop. 9, and the discordance is stated rather than smoothed.
%
% What the folder announces and never establishes: the property (P4), named
% once on p. 15 and never used; the « propriété surjective » of the same
% marginal note, in its sheaf-of-sets form; the definition of « espace
% noethérien spécial » and of Loc_x, announced on p. 15 in a line that is
% almost entirely illegible; the statements 2.7, 2.8, 2.9 and 2.11; and the
% Bourbaki exposé number of 2.17, left as four dots on the leaf.

\documentclass[11pt,a4paper]{article}
% Le préambule commun à toutes les transcriptions du fonds Grothendieck.
%
% Il tient en une idée : **l'incertitude se compose, elle ne se résout pas.**
% Une transcription qui lisse un mot douteux a détruit la seule chose pour
% laquelle on la faisait — un lecteur doit pouvoir savoir, à chaque endroit,
% ce qui était sur le feuillet et ce qui a été ajouté. Les six macros qui
% suivent sont donc l'appareil critique, et non de la décoration.
%
% Les mêmes commandes sont reconnues par `scripts/render.mjs`, qui produit la
% vue de lecture du site. Une macro ajoutée ici doit l'être aussi là-bas,
% faute de quoi le rendu échouera bruyamment — ce qui est voulu : un rendu
% silencieusement faux est pire qu'un rendu absent.

\ProvidesPackage{grothendieck}[2026/08/08 Transcriptions du fonds Grothendieck]

\usepackage{amsmath,amssymb,amsthm}
\usepackage{mathrsfs}
\usepackage{stmaryrd}
% Les diagrammes commutatifs sont partout dans ce fonds : c'est la langue
% même de la géométrie algébrique de Grothendieck, et une transcription qui
% les rendrait en prose ne transcrirait rien.
\usepackage{tikz-cd}
% Français par défaut — c'est la langue des feuillets — mais l'anglais est
% chargé aussi : la traduction et le résumé sont des documents anglais, et la
% typographie française (espaces avant les deux-points) y serait une faute.
% Ces documents basculent par \selectlanguage{english} en tête de corps.
\usepackage[english,french]{babel}
\usepackage{fontspec}
\usepackage{geometry}
\geometry{a4paper,margin=25mm,marginparwidth=28mm}
\usepackage{marginnote}
\usepackage{xcolor}
% ulem, not soul. `soul`'s \st and \ul are built on a character-by-character
% scan that XeLaTeX's font machinery breaks: the document fails with an
% undefined control sequence rather than degrading. ulem does the same two jobs
% and is Unicode-engine safe. `normalem` stops it hijacking \emph, which the
% transcriptions use for Grothendieck's own emphasis.
\usepackage[normalem]{ulem}
\usepackage{hyperref}
\hypersetup{colorlinks=true,linkcolor=black,urlcolor=black,pdfborder={0 0 0}}

% --- Métadonnées du lot -----------------------------------------------------
% Reprises telles quelles par le script de rendu, qui les affiche en tête de
% la vue de lecture. Elles fixent ce que le fichier prétend couvrir : sans
% elles, une transcription flotte sans point d'attache dans un fonds de seize
% mille feuillets.

\newcommand{\folder}[1]{\gdef\grothFolder{#1}}
\newcommand{\batch}[1]{\gdef\grothBatch{#1}}
\newcommand{\leaves}[2]{\gdef\grothFirst{#1}\gdef\grothLast{#2}}
\newcommand{\foldertitle}[1]{\gdef\grothFoldertitle{#1}}
\gdef\grothFolder{?}\gdef\grothBatch{?}\gdef\grothFirst{?}\gdef\grothLast{?}\gdef\grothFoldertitle{}

% --- Le résumé --------------------------------------------------------------
%
% Il ouvre la lecture modernisée, sous le titre « Résumé », et il est composé
% en petit. Ce n'est pas un ornement : c'est le seul passage du document qui
% ne soit pas des mathématiques, et un lecteur doit pouvoir le reconnaître
% avant de l'avoir lu — puis le sauter s'il connaît déjà le sujet, ou s'y
% arrêter s'il ne le connaît pas. Le filet qui le suit marque la reprise du
% corps.

\newenvironment{resume}
  {\small\setlength{\parskip}{0.45em}}
  {\par\medskip\hrule\medskip}

% --- Mots-clés ---------------------------------------------------------------
%
% En anglais, en clôture du résumé de la lecture modernisée : le vocabulaire
% moderne sous lequel chercher ce que les feuillets construisent. Le manifeste
% du site (`npm run manifest`) les extrait pour étiqueter la cote entière —
% cette ligne est la source unique des tags, il n'y a pas de fichier à côté.
\newcommand{\keywords}[1]{\par\smallskip\noindent
  {\footnotesize\textsc{Keywords} --- \textit{#1}}\par}

% --- L'appareil critique ----------------------------------------------------

% Le numéro de feuillet, en marge. C'est l'ancre qui permet de régler un
% désaccord sur une formule en regardant *un* feuillet plutôt qu'une chemise
% de six cents. Le site s'en sert aussi pour tourner les pages du fac-similé
% quand on fait défiler la transcription.
\newcommand{\leaf}[1]{%
  \marginnote{\footnotesize\color{gray!70}\textsf{#1}}%
  \label{leaf:#1}\ignorespaces}

% Illisible : on ne devine pas. Un mot inventé qui se lit comme les autres est
% le pire résultat possible.
\newcommand{\ill}{\textcolor{red!60!black}{[\,\ldots\,]}}

% Lecture douteuse : le mot est donné, et signalé comme incertain.
\newcommand{\uncertain}[1]{\uline{#1}}

% Ajout de l'éditeur — une lettre manquante, un mot sous-entendu.
\newcommand{\add}[1]{\textcolor{blue!50!black}{[#1]}}

% Biffé par Grothendieck lui-même. Ce qu'il a rayé fait partie du document :
% une rature dit souvent où la pensée a changé de direction.
\newcommand{\struck}[1]{\sout{#1}}

% Note du transcripteur, sur l'état matériel du feuillet.
\newcommand{\note}[1]{\par\smallskip{\small\color{gray!60!black}\textit{[#1]}}\par\smallskip}

% Note marginale de Grothendieck, distincte du corps du texte.
\newcommand{\marginal}[1]{\par\smallskip\noindent\hspace{1em}%
  {\small\color{gray!60!black}#1}\par\smallskip}

% --- Titre ------------------------------------------------------------------

\AtBeginDocument{%
  \begin{center}
    {\large\bfseries Fonds Alexandre Grothendieck --- cote n\textsuperscript{o}~\grothFolder}\\[2pt]
    {\itshape\grothFoldertitle}\\[4pt]
    {\small feuillets \grothFirst--\grothLast\ (lot \grothBatch)}\\[2pt]
    {\footnotesize\color{gray!60!black}Transcription établie d'après le fac-similé numérisé
     par l'Université de Montpellier.\\
     Les lectures incertaines sont soulignées, les illisibles notées [\,\ldots\,],
     les ajouts entre crochets.}
  \end{center}
  \vspace{1em}}

\endinput


\folder{21}
\pages{1}{28}
\dating{[vers 1960-1967]}
\watermark{Édition de démonstration\\interprétation personnelle de l'œuvre}
\foldertitle{[Chapitre] 0 III : tapuscrit et copies de tapuscrit annotés (s.d.), notes manuscrites (s.d.). — lecture modernisée du dossier entier}

\begin{document}

\section*{Résumé}

\begin{resume}
Ce dossier pose une question qu'on peut énoncer sans aucun vocabulaire
technique : \emph{si je connais un objet partout sauf sur une petite partie
fermée, est-ce que je le connais partout ?}

Deux exemples classiques donnent les deux réponses possibles. Prenons le plan
complexe et retirons-en l'origine. La fonction $1/z$ est parfaitement
holomorphe sur ce qui reste, et elle ne se prolonge pas au point retiré : ici,
connaître la fonction en dehors de l'origine ne suffit pas. Prenons maintenant
l'espace complexe de dimension deux et retirons-en encore l'origine. Cette
fois, toute fonction holomorphe sur ce qui reste se prolonge, et de façon
unique, à l'espace entier : c'est un théorème de Hartogs, et il dit que la
réponse a changé. Qu'est-ce qui a changé ? Pas la nature de l'objet retiré — un
point dans les deux cas — mais son \emph{épaisseur relative} : dans le plan
complexe, un point est de codimension un, et l'on peut tourner autour ; dans
l'espace de dimension deux, il est de codimension deux, et il n'y a pas assez
de place autour de lui pour qu'une fonction y prenne un comportement sauvage.

Tout le dossier est une machine à mesurer cette épaisseur, et à en tirer des
conséquences. La mesure s'appelle aujourd'hui la \emph{profondeur} le long de
la partie fermée, et l'on trouve son nom sur la chemise de kraft qui enveloppe
la seconde moitié des feuillets : « profondeur $\geqslant 2$ pour faisceaux
d'ensembles ». Deux conditions sont isolées. La première, notée $(P_1)$, dit
qu'une section est \emph{déterminée} par ce qu'elle vaut en dehors du fermé :
deux sections qui coïncident dehors coïncident partout. La seconde, notée
$(P_2)$, ajoute que toute section définie dehors \emph{provient} d'une section
définie partout. La première est la profondeur $\geqslant 1$, la seconde la
profondeur $\geqslant 2$, et le théorème de Hartogs est le cas où la seconde
est vérifiée.

Le problème, et c'est lui qui occupe les vingt-quatre feuillets, est que ces
deux conditions ne se laissent pas vérifier directement. On voudrait les
ramener à quelque chose de local — les tester point par point — et à quelque
chose de calculable. Le dossier procède en deux temps, et les deux temps sont
matériellement séparés : onze feuillets tapés à la machine, puis treize
feuillets manuscrits.

Les feuillets tapés construisent l'outillage, et ils le font dans le cadre le
plus confortable, celui où l'on peut additionner : les modules sur un anneau.
Deux idées y sont mises en place. La première est un principe d'économie : pour
vérifier qu'une certaine mesure est nulle, il n'est pas nécessaire de la
calculer sur tous les objets, il suffit de la calculer sur une famille
d'objets assez riche pour engendrer tous les autres — exactement comme on
vérifie qu'une application linéaire est nulle en la testant sur une base. La
seconde est un principe de représentation : une mesure qui se comporte bien
vis-à-vis des suites exactes n'est pas un objet informe, c'est toujours
« prendre les homomorphismes vers un seul module fixe » ; et ce module fixe se
fabrique par un passage à la limite explicite. Ce module, le tapuscrit ne lui
donne pas de nom ; c'est celui que l'on appelle aujourd'hui la
\emph{cohomologie locale} le long du fermé.

Les feuillets manuscrits appliquent l'outillage, et ils commencent par le
priver de ce qui le rendait confortable. Un faisceau d'ensembles n'a pas
d'addition : on ne peut ni soustraire deux sections ni parler du noyau d'une
application. Les conditions $(P_1)$ et $(P_2)$, elles, gardent un sens, et la
question devient : le principe de localisation survit-il ? La réponse est oui,
et le prix à payer est l'introduction, en chaque point, d'un espace plus petit
— l'ensemble des points dont le point considéré est une spécialisation — sur
lequel la condition se lit de façon purement locale. La démonstration est une
récurrence sur les parties fermées, c'est-à-dire exactement le même mouvement
que celui qui, dans les feuillets tapés, servait à fabriquer les suites de
composition.

Le dernier tiers du dossier fait un pas de plus, et c'est le plus intéressant.
Jusque-là on prolongeait des \emph{sections} ; on va maintenant prolonger des
\emph{objets} — un fibré, un revêtement, un torseur. La différence est qu'entre
deux objets il n'y a plus égalité ou non-égalité, mais des isomorphismes, en
général plusieurs. Une collection d'objets attachée à chaque ouvert, avec des
restrictions, s'appelle une \emph{catégorie fibrée}, et le foncteur de
restriction d'un ouvert à son intersection avec le complémentaire du fermé peut
être fidèle, pleinement fidèle, ou une équivalence. Là où il y avait deux
conditions, il y en a donc trois : $(P_1)$, $(P_2)$, $(P_3)$. Et le dossier
établit que ces trois conditions se lisent sur les trois premiers groupes de
cohomologie locale $H^0_Y$, $H^1_Y$, $H^2_Y$ — la troisième faisant intervenir
un degré de plus que la seconde, exactement comme la seconde en fait intervenir
un de plus que la première. Le dernier feuillet en donne le cas le plus net,
celui des torseurs sous un faisceau de groupes, où $H^0_Y$ et $H^1_Y$
apparaissent tout simplement comme le noyau et le conoyau de l'application dont
on mesurait au départ l'injectivité et la bijectivité.

Reste à dire ce que ces feuillets sont, matériellement, car cela fait partie de
ce qu'on y lit. La chemise porte, de sa main, « Notes pour Dieud. §9 » : ce
sont des notes destinées à Dieudonné, c'est-à-dire de la matière première pour
un chapitre des \emph{Éléments}. Et le tapuscrit est un texte qui se parle à
lui-même de sa propre rédaction. La première ligne du premier feuillet propose
déjà de déplacer le numéro ailleurs. Une note marginale, deux fois répétée, dit
« déployer ». Une phrase avoue qu'« il aurait fallu déployer la formule
également dans 2.4 ». Une autre : « Tout ceci devrait figurer dans un numéro à
part au Chap 0 ». Une autre encore : « Il serait préférable d'inclure 2.13 dans
l'énoncé de 2.12 […] pour éviter des redites ». Et au milieu du dernier
feuillet tapé, une ligne entière est biffée à la machine, qui commençait :
« Si on rejette tout ce numéro (comme il me semble raisonnable) ». Il est en
train d'écrire, et simultanément de plaider contre l'existence de ce qu'il
écrit. On trouve même, au détour d'un corollaire, la remarque qu'il aurait été
« canularesque » d'énoncer la proposition sous sa forme contravariante. La
seconde moitié est l'exact opposé : une écriture rapide où les énoncés et les
formules passent et où les phrases qui les relient, très souvent, ne passent
pas. Ce sont deux régimes de travail, et le dossier les a conservés côte à
côte.

Les noms sous lesquels on cherchera la suite : cohomologie locale et
profondeur, prolongement à travers un fermé, théorème d'Eilenberg--Watts,
théorie de Matlis des modules injectifs, espaces spectraux, champs et descente,
et — pour le dernier feuillet — cohomologie non abélienne.

\keywords{local cohomology, depth along a closed subset, extension across a closed subset, Hartogs phenomenon, cohomological functor, system of generators, noetherian object, Eilenberg-Watts theorem, tensoring an abelian category over a ring, finitely presented module, Matlis theory, injective hull, associated primes, embedded prime, torsion functor, Artin-Rees lemma, ind-representable functor, coherent sheaf with support in a closed set, sheaf of sets, spectral space, generization, localisation of a space at a point, fibered category, stack, pseudo-limit, 2-colimit, faithful functor, fully faithful functor, noetherian induction, constructible set, torsor, principal homogeneous space, nonabelian cohomology, Lefschetz condition}
\end{resume}

\pagerange{3}{28}
\section*{Le fil du dossier, et les conventions}

\emph{Les stations.}

\begin{enumerate}
  \item \emph{Pages 3 à 5}, tapuscrit : tester l'annulation d'un foncteur
  cohomologique sur un système de générateurs. Propositions 2.1 à 2.5.
  \item \emph{Pages 5 à 7}, tapuscrit : le produit tensoriel $P \otimes_A M$
  d'un module par un objet d'une catégorie abélienne, et la caractérisation des
  foncteurs exacts à droite par leur valeur en $A$. Proposition 2.6.
  \item \emph{Pages 8 à 10}, tapuscrit : le cas contravariant, le passage à la
  limite sur $A_n = A/I^{n+1}$, et le module $H$ qui représente un foncteur
  exact à gauche. Les énoncés 2.7 à 2.11, dont aucun n'est intégralement sur
  les feuillets.
  \item \emph{Pages 10 et 11}, tapuscrit : le critère d'annulation en termes de
  $H$ (2.12, 2.13), et les idéaux premiers associés de $H$ (2.14).
  \item \emph{Pages 11 à 13}, tapuscrit : $H$ injectif si et seulement si le
  foncteur est exact (2.15), la stabilité des injectifs par torsion (2.16), et
  la généralisation aux faisceaux cohérents (2.17).
  \item \emph{Pages 15 à 20}, manuscrit : les propriétés $(P_1)$ et $(P_2)$
  pour un faisceau d'ensembles, la localisation en un point, et l'équivalence
  entre la forme globale et la forme ponctuelle. Définitions 1 et 2,
  Proposition 3, Corollaire 4, Lemme 5.
  \item \emph{Pages 21 à 24}, manuscrit : les trois conditions $(P_1)$,
  $(P_2)$, $(P_3)$ pour une catégorie fibrée, leur traduction en termes des
  faisceaux $\mathcal{H}om(f,g)$, la catégorie fibrée induite en un point, et
  l'énoncé global-ponctuel. Définition 6, Proposition 7, Corollaire 8.
  \item \emph{Pages 25 à 27}, manuscrit : les démonstrations, la récurrence
  noethérienne, et la traduction des trois conditions en l'annulation des
  faisceaux de cohomologie locale. Corollaire 10.
  \item \emph{Page 28}, manuscrit : le cas des torseurs sous un faisceau de
  groupes. Proposition 9.
\end{enumerate}

\emph{La numérotation est la sienne, et elle est en désordre ; elle n'est pas
remise en ordre ici.} Le tapuscrit va de 2.1 à 2.17 sans rupture, mais quatre
de ses numéros — 2.7, 2.8, 2.9, 2.11 — ne portent sur les feuillets aucun
énoncé qu'on puisse lire : 2.9 n'a que son intitulé, 2.10 est coupé par la fin
du feuillet 9, et 2.7, 2.8, 2.11 ne sont connus que par les renvois que le
texte leur fait. Le manuscrit, lui, se numérote Définition 6 (page 21),
Proposition 7 et Corollaire 8 (page 24), Corollaire 10 (page 27), puis
Proposition 9 (page 28) : le 9 vient \emph{après} le 10, sur le dernier
feuillet. Les sections ci-dessous suivent l'ordre des feuillets, non celui des
numéros ; le Corollaire 10 y est donc lu avant la Proposition 9, et c'est
d'ailleurs heureux, car la seconde est ce qui permet de corriger le premier.

\emph{Ce que seule une lecture d'ensemble peut dire.} Trois liens, qu'aucun
feuillet n'énonce.

Le premier, et c'est la charnière du dossier. Le tapuscrit fabrique, à la page
9, un objet
\[
  H^{i} \;=\; \varinjlim_{m} T^{i}(A_{m}) , \qquad A_{m} = A/I^{m+1} ,
\]
et il en fait le représentant du foncteur exact à gauche $T$. Il ne lui donne
aucun nom et n'écrit jamais de cohomologie. Le manuscrit, lui, écrit à la page
27 des faisceaux $\underline{H}^{i}_{Y}(F)$ sans dire d'où ils viennent. Ce
sont le même objet : pour $T^{i}(M) = \mathrm{Ext}^{i}(M,F)$, la limite
inductive des $T^{i}(A_{m})$ est la cohomologie locale de $F$ le long de
$Y = V(I)$. Le tapuscrit est la \emph{construction} de l'objet dans lequel le
manuscrit \emph{énonce} son critère. C'est la seule raison pour laquelle les
deux paquets sont dans la même chemise, et ni l'un ni l'autre ne le dit.

Le deuxième : la récurrence noethérienne est un seul mouvement, employé deux
fois dans deux catégories différentes. À la page 3, pour montrer qu'un objet
admet une suite de composition, on prend un sous-objet \emph{maximal} qui en
admet une et l'on montre qu'il est le tout ; aux pages 19 et 26, pour établir
une propriété relativement à un fermé $Y$, on la suppose acquise pour tout
fermé strictement plus petit et l'on descend en un point \emph{maximal} de $Y$.
Le tapuscrit démontre la forme abstraite de ce que le manuscrit utilise sous
forme topologique.

Le troisième : la note marginale de la page 15, « N.B. regarder aussi la
propriété : \emph{surjective} », reste sans réponse dans la partie consacrée
aux faisceaux d'ensembles. Elle est pourtant la question que les pages 24 à 27
traitent exclusivement. Car pour une catégorie fibrée, la surjectivité
essentielle — tout objet au-dessus de $V \cap U$ est la restriction d'un objet
au-dessus de $V$ — est précisément ce que la condition $(P_3)$ ajoute à
$(P_2)$, et c'est le seul des trois cas dont la démonstration coûte quelque
chose.

\emph{Les notations, valables pour tout le dossier.} $A$ est un anneau
commutatif noethérien, $I$ un idéal, $A_n = A/I^{n+1}$, $X = \mathrm{Spec}\, A$,
$Y = V(I)$ la partie fermée qu'il définit, $U = X - Y$ son complémentaire et
$i : U \to X$ l'inclusion. Le complémentaire est noté $U$ partout : le
tapuscrit ne le nomme pas, les pages 15 à 20 l'écrivent $U$ et les pages 21 à
28 l'écrivent $\mathcal{U}$. Les objets d'une catégorie abélienne sont notés
$M$, $N$ ; la lettre $F$ est réservée aux faisceaux et $\mathcal{F}$ aux
catégories fibrées, alors que le tapuscrit écrit $F$ pour un objet quelconque
de $\underline{C}$. Ses catégories $\underline{C}$, $\underline{C}'$,
$\underline{G}$ sont écrites $\mathcal{C}$, $\mathcal{C}'$, $\mathcal{G}$. Pour
un point $x$ d'un espace $X$, $X_{(x)}$ désigne l'ensemble des générisations de
$x$ — les points dont $x$ est une spécialisation — muni de la topologie
induite, ce qu'il note $\mathrm{Loc}_x(X)$, et $X'_{(x)} = X_{(x)} - \{x\}$. Le
faisceau des homomorphismes est noté $\mathcal{H}om$.

\pagerange{3}{13}
\section*{I. Le tapuscrit : l'outillage additif (pages 3 à 13)}

\pagerange{3}{5}
\subsection*{Tester l'annulation sur un système de générateurs (pages 3 à 5)}

Le tapuscrit s'ouvre sur une remarque de rédaction — le numéro pourrait aussi
bien figurer au chapitre $0_{\mathrm{III}}$, qui contiendrait alors certains des
sorites du paragraphe 7 — puis entre immédiatement dans le principe d'économie
qui gouverne toute la première moitié.\footnote{Le feuillet renvoie à un
« Chap. $\mathrm{O}_{III}$ » et à un « par. 7 » sans autre précision. Rien n'a
été confronté ici aux \emph{Éléments} imprimés, conformément à la
transcription, qui ne les a pas ouverts non plus.}

\emph{Proposition 2.1.} Soient $\mathcal{C}$ une catégorie abélienne et
$\mathcal{G} \subset \mathrm{Ob}\,\mathcal{C}$ une partie telle que tout objet
de $\mathcal{C}$ admette une suite de composition finie dont les facteurs sont
isomorphes à des quotients d'objets de $\mathcal{G}$. Soit $\mathcal{C}'$ une
seconde catégorie abélienne.

\begin{enumerate}
  \item[(i)] Un foncteur additif exact à droite $T : \mathcal{C} \to
  \mathcal{C}'$ est nul si et seulement s'il est nul sur les objets de
  $\mathcal{G}$.
  \item[(ii)] Soit $(T^{i})$ un foncteur cohomologique de $\mathcal{C}$ dans
  $\mathcal{C}'$ tel que $T^{i} = 0$ pour $i > i_{0}$, et soit $n$ un entier.
  Alors $T^{i} = 0$ pour $i > n$ si et seulement si $T^{i}(M) = 0$ pour tout
  $M \in \mathcal{G}$ et tout $i > n$.
\end{enumerate}

La démonstration de (ii) est une récurrence \emph{descendante} sur $n$, et
c'est sa mécanique qui mérite d'être retenue, car c'est elle que le dossier
réemploiera. Pour $n \geqslant i_{0}$ l'énoncé est vide. Supposons-le acquis au
rang $n+1$ et supposons $T^{i}(M) = 0$ pour $M \in \mathcal{G}$ et $i > n$.
L'hypothèse de récurrence donne alors $T^{i} = 0$ pour $i > n+1$ ; mais un
foncteur cohomologique dont tous les termes de degré supérieur s'annulent est
exact à droite en son plus haut degré non trivial, donc $T^{n+1}$ est exact à
droite, et (i) le rend nul. La suite exacte longue fait tomber les degrés un à
un, et chaque degré est ramené au cas (i).

\emph{Corollaire 2.2.} Sous les mêmes hypothèses, soit $(T^{i})$ un foncteur
cohomologique \emph{contravariant} de $\mathcal{C}$ dans $\mathcal{C}'$ tel que
$T^{i} = 0$ pour $i < i_{0}$, et soit $n$ un entier. Alors $T^{i} = 0$ pour
$i < n$ si et seulement si $T^{i}(M) = 0$ pour tout $M \in \mathcal{G}$ et tout
$i < n$.

C'est 2.1 (ii) appliqué à $\mathcal{C}$ et $\mathcal{C}'^{\,\circ}$ : un
foncteur contravariant de $\mathcal{C}$ dans $\mathcal{C}'$ est un foncteur
covariant de $\mathcal{C}$ dans $\mathcal{C}'^{\,\circ}$, l'exactitude s'y
renverse, et la suite exacte longue y est lue de haut en bas ; en posant
$S^{j} = T^{-j}$ on est exactement dans les hypothèses de 2.1 (ii) avec
$-i_{0}$ et $-n$.\footnote{Le feuillet ne détaille pas ce renversement ; il
écrit seulement « il suffit d'appliquer 2.1. (ii) à $\underline{C}$ et
$\underline{C}'^{o}$ ». Le calcul des indices est ajouté ici, parce que c'est le
seul endroit du dossier où une variance est retournée et qu'une erreur d'un
signe y rendrait faux tout ce qui suit.} Il ajoute, et c'est sa manière : « pour
qu'il n'y ait pas de jaloux, il faudrait aussi répéter (i) » — de fait, l'énoncé
contravariant de (i) n'est nulle part écrit, et le Corollaire 2.3 le citera
pourtant sous la forme « les conclusions de 2.1 (i)(ii) et 2.2 (i)(ii)
».\footnote{2.2 tel qu'il est énoncé n'a pas de partie (i) ; le renvoi de 2.3 en
suppose une. Elle est immédiate — un foncteur contravariant exact à gauche est
nul dès qu'il l'est sur $\mathcal{G}$ — et c'est elle que la lecture suppose
partout où 2.3 invoque « 2.2 (i) ».} Il note aussi qu'« il aurait été
canularesque d'énoncer 2.1 avec des foncteurs contravariants ».

\emph{Corollaire 2.3.} Soient $\mathcal{C}$, $\mathcal{C}'$ deux catégories
abéliennes et $\mathcal{G} \subset \mathrm{Ob}\,\mathcal{C}$. On suppose que les
objets de $\mathcal{C}$ sont \emph{noethériens}, c'est-à-dire satisfont la
condition de chaîne ascendante pour les sous-objets, et que $\mathcal{G}$ est un
\emph{système de générateurs} de $\mathcal{C}$, au sens suivant : pour tout
objet $N \neq 0$ de $\mathcal{C}$, il existe $M \in \mathcal{G}$ et un
homomorphisme non nul $M \to N$.\footnote{Le feuillet écrit « il existe un
$M \in \underline{C}$ », ce qui rendrait la condition vide : $N$ lui-même
convient toujours. C'est $M \in \underline{G}$ qu'il faut, et c'est ce qui est
écrit ici ; toute la démonstration qui suit l'utilise sous cette forme.} Alors
les conclusions de 2.1 et de 2.2 valent.

La démonstration est le premier des deux emplois de la récurrence noethérienne
annoncés plus haut. Soit $N$ un objet de $\mathcal{C}$ et soit $N'$ un
sous-objet \emph{maximal} parmi ceux qui admettent une suite de composition du
type voulu — il en existe, car $0$ en admet une et les objets sont noethériens.
Si $N' \neq N$, alors $N/N' \neq 0$, donc il existe $M \in \mathcal{G}$ et un
homomorphisme non nul $M \to N/N'$ ; l'image inverse dans $N$ de son image est
un sous-objet strictement plus grand que $N'$ qui admet encore une telle suite,
ce qui contredit la maximalité. Donc $N' = N$.

\emph{Exemple 2.4.} Soient $A$ un anneau commutatif noethérien, $I$ un idéal,
$X = \mathrm{Spec}\, A$, $Y = V(I)$. Soit $\mathcal{C}$ la catégorie des
$A$-modules de type fini dont tout élément est annulé par une puissance de $I$
— de façon équivalente, la catégorie des modules cohérents sur $X$ à support
contenu dans $Y$.\footnote{Le feuillet écrit d'abord « Modules quasi-cohérents », la mention
\emph{quasi-} étant ensuite surfrappée à la machine : c'est
bien la catégorie des modules cohérents, comme l'exige l'hypothèse « de type
fini » qu'il ajoute à l'encre sur la même ligne.} Ses objets sont noethériens,
et pour $M \in \mathcal{C}$ :
\[
  M \text{ est un générateur de } \mathcal{C}
  \quad\Longleftrightarrow\quad
  \mathrm{supp}\, M = Y .
\]

La nécessité : si $\mathrm{supp}\, M \neq Y$, une composante irréductible
$Y_{j}$ de $Y$ n'est pas contenue dans $\mathrm{supp}\, M$ ; en prenant
$N = \mathcal{O}_{Y_{j}}$ pour la structure réduite induite, on a
$\mathrm{Hom}(M,N) = 0$ avec $N \neq 0$. La suffisance repose sur la formule
\[
  \mathrm{Ass}\,\mathcal{H}om(M,N) \;=\; \mathrm{supp}\, M \cap \mathrm{Ass}\, N ,
\]
valable pour $M$ de type fini : si $\mathrm{supp}\, M = Y$ et si $N$ est
quasi-cohérent de support contenu dans $Y$, alors $\mathrm{Ass}\, N \subset Y =
\mathrm{supp}\, M$, de sorte que $\mathrm{Hom}(M,N) = 0$ force
$\mathrm{Ass}\, N = \emptyset$, donc $N = 0$. Le même argument montre qu'une
famille $(M_{j})$ d'objets de $\mathcal{C}$ est un système de générateurs si et
seulement si
\[
  \bigcup_{j} \mathrm{supp}\, M_{j} \;=\; Y ,
\]
et donc — c'est le point — que pour vérifier l'annulation d'un foncteur comme
dans 2.1 ou 2.2, il suffit de le tester sur les $M_{j}$, en particulier sur un
seul module de support $Y$.

\emph{Exemple 2.5.} La même chose sur un schéma : $X$ noethérien
quasi-compact muni d'un faisceau ample $\mathcal{O}_{X}(1)$, $Y$ une partie
fermée, $\mathcal{C}$ la catégorie des faisceaux cohérents à support contenu
dans $Y$. Si $(M_{j})$ est une famille d'objets de $\mathcal{C}$ telle que
$\bigcup_{j} \mathrm{supp}\, M_{j} = Y$, alors les $M_{j}(-n)$ pour $n \geqslant
N$ forment un système de générateurs de $\mathcal{C}$, et même de la catégorie
de tous les modules quasi-cohérents de support contenu dans
$Y$.\footnote{L'énoncé est écrit sur le feuillet avec $n > N$ et sa conclusion
avec $n \geqslant N$ ; l'un des deux est un lapsus et la différence est sans
conséquence, la famille étant cofinale dans les deux cas. $n \geqslant N$ est
retenu ici.} En effet, si $F$ est quasi-cohérent non nul de support contenu dans
$Y$, on se place sur un ouvert affine où $F \neq 0$ et 2.4 donne
$\mathcal{H}om(M_{j},F) \neq 0$ pour un $j$ convenable ; l'amplitude fournit
alors des sections non nulles de $\mathcal{H}om(M_{j},F)(n)$ pour $n$ grand,
c'est-à-dire des homomorphismes non nuls $M_{j}(-n) \to F$.

\pagerange{5}{7}
\subsection*{Le produit tensoriel $P \otimes_A M$, et les foncteurs exacts à droite (pages 5 à 7)}

Le tapuscrit change alors d'objet. $A$ est maintenant un anneau quelconque, non
nécessairement commutatif, $\mathcal{C}$ une catégorie abélienne de
$A$-modules à gauche — par exemple n'importe quelle sous-catégorie pleine de la
catégorie des $A$-modules stable par noyaux et conoyaux — et $\mathcal{C}'$ une
catégorie abélienne.\footnote{Le feuillet dit bien « une catégorie abélienne de
$A$-modules », et non « la catégorie » ; la transcription le signale, et la
distinction compte, puisque tout ce qui suit vaut pour une sous-catégorie
pleine stable par noyaux et conoyaux.}

Soit $P$ un objet de $\mathcal{C}'$ muni d'une structure de $A$-module à droite,
c'est-à-dire d'un homomorphisme d'anneaux $A^{\circ} \to \mathrm{End}(P)$. Pour
tout $Q \in \mathcal{C}'$, le groupe $\mathrm{Hom}_{\mathcal{C}'}(P,Q)$ est
alors un $A$-module à \emph{gauche}, par $a \cdot u = u \circ \rho_{a}$. On pose,
lorsqu'il existe, $P \otimes_{A} M$ pour l'objet de $\mathcal{C}'$ défini par la
propriété universelle
\[
  \mathrm{Hom}_{\mathcal{C}'}(P \otimes_{A} M,\, Q)
  \;=\;
  \mathrm{Hom}_{A}(M,\, \mathrm{Hom}_{\mathcal{C}'}(P,Q)) ,
\]
fonctorielle en $Q$.

Cet objet n'existe pas toujours, et le feuillet énumère exactement ce qui le
fait exister. D'abord $P \otimes_{A} A_{g} \simeq P$, où $A_{g}$ est $A$ vu
comme $A$-module à gauche. Ensuite, si $M'' \to M' \to M \to 0$ est exacte et si
$P \otimes_{A} M'$ et $P \otimes_{A} M''$ existent, alors $P \otimes_{A} M$
existe et
\[
  P \otimes_{A} M'' \longrightarrow P \otimes_{A} M' \longrightarrow
  P \otimes_{A} M \longrightarrow 0
\]
est exacte : sous réserve d'existence, $P \otimes_{A} M$ est exact à droite en
$M$. Enfin le foncteur commute aux limites inductives quelconques — pas
seulement filtrantes — au sens précis suivant : si les $P \otimes_{A} M_{j}$
existent, alors $P \otimes_{A} \varinjlim M_{j}$ existe si et seulement si
$\varinjlim (P \otimes_{A} M_{j})$ existe, et les deux sont canoniquement
isomorphes. En combinant : $P \otimes_{A} M$ existe dès que $M$ est de
présentation finie, car il existe pour $A_{g}^{(I)}$ avec $I$ fini, donc pour un
conoyau $A_{g}^{(I)} \to A_{g}^{(J)}$ ; et il existe sans restriction sur $M$ si
$\mathcal{C}'$ admet les sommes directes quelconques. La construction se lit sur
une présentation : la matrice qui présente $M$ définit un morphisme
$P^{(I)} \to P^{(J)}$ dont le conoyau est l'objet cherché.\footnote{« Tout ceci
devrait figurer dans un numéro à part au Chap 0, sans trainer l'accent $'$ sur
la catégorie abélienne », note-t-il ici ; et en marge, deux fois dans le
tapuscrit, le seul mot « déployer ».}

La version duale est obtenue en remplaçant $\mathcal{C}'$ par l'opposée d'une
catégorie abélienne $\mathcal{C}_{1}$ : les $A$-$\mathcal{C}'$-modules à droite
sont alors les $A$-$\mathcal{C}_{1}$-modules à gauche, et le produit tensoriel
s'écrit, en termes de $\mathcal{C}_{1}$, $\mathrm{Hom}_{A}(M,P)$, défini
directement par
\[
  \mathrm{Hom}_{\mathcal{C}_{1}}(Q,\, \mathrm{Hom}_{A}(M,P))
  \;=\;
  \mathrm{Hom}_{A}(M,\, \mathrm{Hom}_{\mathcal{C}_{1}}(Q,P)) .
\]
« Il y aurait donc lieu de développer les deux formalismes en parallèle », l'un
sur les produits tensoriels, l'autre sur les Hom.\footnote{La ligne qui précède
ce dernier déploiement n'a pas été lue sur le feuillet : la transcription
signale que les deux membres de la formule sont les seuls à y figurer en clair
et que la phrase qui les introduit est reconstituée du contexte. La formule
elle-même ne fait aucun doute, étant la définition par adjonction.}

Soit alors $\mathcal{C}$ la catégorie des $A$-modules à gauche de présentation
finie. Tout $A$-$\mathcal{C}'$-module à droite $P$ donne un foncteur additif
$M \rightsquigarrow P \otimes_{A} M$ de $\mathcal{C}$ dans $\mathcal{C}'$, et
inversement :

\emph{Proposition 2.6.} Soit $T : \mathcal{C} \to \mathcal{C}'$ un foncteur
additif, et soit $P = T(A)$, muni de sa structure naturelle de $A$-module à
droite. L'homomorphisme canonique $P \otimes_{A} M \to T(M)$ est un isomorphisme
si et seulement si $T$ est exact à droite.

C'est le théorème qu'on appelle aujourd'hui d'Eilenberg--Watts, ici sous la
forme où la catégorie but est une catégorie abélienne quelconque et non la
catégorie des modules sur un anneau.\footnote{Les deux articles fondateurs —
Eilenberg, \emph{Abstract description of some basic functors}, et Watts,
\emph{Intrinsic characterizations of some additive functors} — datent de 1960,
et le dossier est daté par les archivistes « vers 1960-1967 ». Le feuillet ne
cite personne, ni ne nomme le résultat. Rien ici n'autorise à dire qu'il les
connaissait, ni le contraire.}

Le feuillet 7 n'a pu être lu de bout en bout ; il porte encore le numéro 2.7,
dont l'énoncé n'est pas accessible.\footnote{La transcription le dit
explicitement : « la page 7 n'a pu être lue de bout en bout ; le texte ci-dessus
en donne les énoncés, les phrases de liaison manquent ». 2.7 est connu par le
seul renvoi de la page 8, « il faudra expliciter les énoncés duals de 2.6. et
2.7. », qui le place à côté de 2.6 et suggère qu'il s'agit de la forme
catégorique de celle-ci — l'équivalence entre foncteurs additifs exacts à droite
$\mathcal{C} \to \mathcal{C}'$ et $A$-$\mathcal{C}'$-modules à droite. Cette
identification est une inférence, et n'est pas sur le feuillet.} Un
homomorphisme canonique $P \otimes_{A} M \to \,?$ y est amorcé dont le second
membre n'a pas été lu.

\pagerange{8}{10}
\subsection*{Le cas contravariant, et le module qui représente (pages 8 à 10)}

La page 8 reprend en cours de phrase, sur le cas contravariant : un foncteur
$T : \mathcal{C}^{\circ} \to \mathcal{C}'$, ce qui se ramène au précédent en
remplaçant $\mathcal{C}'$ par $\mathcal{C}'^{\,\circ}$. La contravariance
retourne la structure de module : $T(A)$ porte maintenant une structure de
$A$-module à \emph{gauche}, et l'homomorphisme canonique, fonctoriel en $M$, est
\[
  T(M) \longrightarrow \mathrm{Hom}_{A}(M, T(A)) .
\]
C'est un isomorphisme si et seulement si $T$ est \emph{exact à gauche},
c'est-à-dire transforme une suite exacte $M'' \to M' \to M \to 0$ de
$\mathcal{C}$ en une suite exacte $0 \to T(M) \to T(M') \to T(M'')$. On obtient
ainsi une équivalence entre la catégorie des foncteurs contravariants exacts à
gauche de $\mathcal{C}$ dans $\mathcal{C}'$ et la catégorie, notée $A(\mathcal{C}')$,
des objets de $\mathcal{C}'$ munis d'une structure de $A$-module à
gauche.\footnote{Le retournement du côté est ce que la contravariance impose et
n'est pas dit sur le feuillet : dans le cas covariant $T(A)$ est un module à
droite, dans le cas contravariant un module à gauche. La notation
$A(\underline{C}')$ du feuillet, où le $A$ est ajouté à l'encre devant la
parenthèse, est cohérente avec cette lecture, puisqu'il écrit ailleurs
« $A$-$\underline{C}'$-modules à droite » lorsqu'il veut l'autre côté.} Le
résultat s'étend à la catégorie de tous les $A$-modules, sans hypothèse
noethérienne sur $A$, pourvu que les produits quelconques existent dans
$\mathcal{C}'$.

Vient alors le passage qui donne au tapuscrit sa raison d'être. On revient à la
catégorie $\mathcal{C}$ de l'exemple 2.4 : les $A$-modules de type fini annulés
par une puissance de $I$. Pour $n \geqslant 0$, soit $\mathcal{C}_{n}$ la
sous-catégorie pleine des modules de type fini sur $A$ annulés par $I^{n+1}$,
équivalente à la catégorie des modules de type fini sur $A_{n} = A/I^{n+1}$. Le
foncteur $\mathcal{C}_{n} \to \mathcal{C}$ est exact et $\mathcal{C}$ est
réunion filtrante des $\mathcal{C}_{n}$.\footnote{Le feuillet écrit « les modules
de type fini sur $\underline{C}$ » là où l'argument demande « sur $A$ » ; la
transcription porte la variante telle quelle et la signale. La lecture retenue
ici est « sur $A$ », la seule qui donne un sens à la phrase suivante, qui
identifie $\mathcal{C}_{n}$ aux modules de type fini sur $A_{n}$. Le même
feuillet écrit d'abord $\underline{C}^{n}$ puis $\underline{C}_{n}$ pour la même
sous-catégorie ; l'indice inférieur est retenu.}

Un foncteur additif $T : \mathcal{C} \to \mathcal{C}'$ est déterminé par ses
restrictions $T_{n}$ aux $\mathcal{C}_{n}$, et ce qui précède fournit pour
chacune un homomorphisme fonctoriel $T(A_{n}) \otimes_{A} M \to T_{n}(M)$. Pour
$M$ fixé dans un $\mathcal{C}_{n_{0}}$, la famille $(T(A_{n}) \otimes_{A} M)_{n
\geqslant n_{0}}$ est un système projectif, dont le système projectif constant
de valeur $T(M)$ est un quotient ; sous réserve d'existence des limites, on
obtient
\[
  \left(\varprojlim_{n} T(A_{n})\right) \otimes_{A} M \longrightarrow T(M) ,
\]
et cet homomorphisme ne dépend pas du choix de $n_{0}$ tel que
$M \in \mathcal{C}_{n_{0}}$.

C'est cependant la situation \emph{duale} dont le dossier aura besoin. Pour
$T : \mathcal{C}^{\circ} \to \mathcal{C}'$ contravariant, on a des
homomorphismes
\[
  (\ast) \qquad
  T_{n}(M) \longrightarrow \mathrm{Hom}_{A}(M, T_{n}(A_{n}))
  \qquad (n \geqslant n_{0}) ,
\]
compatibles au passage à la limite, d'où
\[
  T(M) \longrightarrow \varinjlim_{n} \mathrm{Hom}_{A}(M, T(A_{n}))
  \longrightarrow \mathrm{Hom}_{A}\!\left(M, \varinjlim_{n} T(A_{n})\right) .
\]
Pour que $T$ soit exact à gauche, il faut et il suffit que les $T_{n}$ le
soient, ou encore que les $(\ast)$ soient des isomorphismes. Et si dans
$\mathcal{C}'$ les limites inductives filtrantes dénombrables existent et sont
exactes — a fortiori si l'axiome (AB 5) est vérifié — le foncteur
$P \rightsquigarrow \mathrm{Hom}_{A}(M,P)$ commute à ces limites pour tout $M$
de présentation finie sur $A$, de sorte que l'homomorphisme composé
\[
  T(M) \longrightarrow \mathrm{Hom}_{A}\!\left(M, \varinjlim_{n} T(A_{n})\right)
\]
est un isomorphisme. La réciproque est triviale.

\emph{Proposition 2.9.} Soient $A$ un anneau commutatif noethérien, $I$ un
idéal, $\mathcal{C}$ la catégorie des $A$-modules de type fini annulés par une
puissance de $I$, et $\mathcal{C}'$ une catégorie abélienne où les limites
inductives filtrantes dénombrables existent et sont exactes. Pour un foncteur
contravariant $T : \mathcal{C}^{\circ} \to \mathcal{C}'$, les conditions
suivantes sont équivalentes : $T$ est exact à gauche ; l'homomorphisme canonique
$T(M) \to \mathrm{Hom}_{A}(M, H)$ est un isomorphisme pour tout
$M \in \mathcal{C}$, où l'on a posé
\[
  H \;=\; \varinjlim_{n} T(A_{n}) .
\]

\footnote{Sur le feuillet 9, 2.9 n'est qu'un intitulé : « \emph{Proposition
2.9.} », suivi d'un blanc, puis le texte reprend par « généralisant 2.8 sous les
conditions dites ». L'énoncé ci-dessus n'est pas de lui : c'est celui que le
paragraphe qui le précède immédiatement vient d'établir mot pour mot, et il est
reconstitué à ce titre. Rien n'interdit que ce qu'il avait en tête fût plus
général.}

La suite du feuillet 9 est perdue avec son bas, et la page 10 s'ouvre en cours
de phrase. Ce qu'on en lit décrit la catégorie qui apparaît dans le
\emph{Corollaire 2.10} : les systèmes inductifs $(H_{n})$ dont chaque morphisme
de transition $H_{n} \to H_{m}$ induit un isomorphisme de $H_{n}$ sur le
sous-objet $\mathrm{Hom}_{A}(A_{n}, H_{m})$ des « éléments annulés par
$I^{n+1}$ ». Sous les hypothèses de 2.9, cette catégorie est équivalente à la
catégorie des $A$-$\mathcal{C}'$-modules dont 2.10 parle, laquelle est donc un
cas particulier d'un énoncé 2.11.\footnote{L'énoncé de 2.10 est coupé net par la
fin du feuillet 9 — il s'arrête sur « la catégorie des » — et 2.11 tombe
entièrement dans la lacune entre les deux feuillets : il n'est connu que par le
renvoi « qui est donc un cas particulier de 2.11 ». Ni l'un ni l'autre n'est
reconstitué ici. Ce qui est sûr, et suffit à la suite du dossier, est
l'équivalence entre foncteurs contravariants exacts à gauche sur $\mathcal{C}$ et
modules de $I$-torsion, par $T \mapsto \varinjlim T(A_{n})$ et
$H \mapsto \mathrm{Hom}_{A}(-,H)$.}

\pagerange{10}{11}
\subsection*{Le critère d'annulation, et les premiers associés (pages 10 et 11)}

\emph{Corollaire 2.12.} Soient $A$ un anneau commutatif noethérien, $I$ un
idéal, $Y = V(I)$, $\mathcal{C}$ la catégorie des $A$-modules de type fini
annulés par une puissance de $I$, et $\mathcal{C}'$ une catégorie abélienne où
les limites inductives filtrantes dénombrables sont exactes. Soit $(T^{i})$ un
foncteur cohomologique contravariant de $\mathcal{C}$ dans $\mathcal{C}'$ tel
que $T^{i} = 0$ pour $i < i_{0}$, soit $n$ un entier, et soit $(M_{j})$ une
famille de $A$-modules de type fini telle que $\bigcup_{j} \mathrm{supp}\, M_{j}
= Y$. Posons
\[
  H^{i} \;=\; \varinjlim_{m} T^{i}(A_{m}) .
\]
Les conditions suivantes sont équivalentes :
\begin{enumerate}
  \item[(i)] $T^{i} = 0$ pour $i < n$ ;
  \item[(ii)] $H^{i} = 0$ pour $i < n$ ;
  \item[(iii)] $T^{i}(M_{j}) = 0$ pour $i < n$ et tout $j$.
\end{enumerate}

L'équivalence de (i) et (iii) est 2.2 combinée à 2.4 : la famille $(M_{j})$ est
un système de générateurs. L'implication (i) $\Rightarrow$ (ii) est immédiate.
Pour (ii) $\Rightarrow$ (i) on procède par récurrence sur $n$ ; c'est vide pour
$n \leqslant i_{0}$, et supposant l'énoncé acquis pour tous les $n' < n$, il
reste à voir que $T^{n-1} = 0$. Or les degrés inférieurs étant nuls, la suite
exacte longue rend $T^{n-1}$ exact à gauche, donc 2.9 donne $T^{n-1}(M) \simeq
\mathrm{Hom}_{A}(M, H^{n-1})$, qui est nul puisque
$H^{n-1} = 0$.\footnote{Le mot ajouté à l'encre au-dessus de cette ligne se lit
« nul » ou « une » ; la transcription le signale comme douteux et retient
« nul », qui est ce que l'argument demande. La lecture est ici la même.}

\emph{Corollaire 2.13.} Sous ces conditions, on a un isomorphisme fonctoriel
\[
  T^{n}(M) \;=\; \mathrm{Hom}_{A}(M, H^{n}) .
\]
\emph{Il serait préférable}, note-t-il aussitôt, \emph{d'inclure 2.13 dans
l'énoncé de 2.12 et de démontrer par récurrence} (ii) $\Rightarrow$ (i)
\emph{plus l'énoncé final 2.13, pour éviter des redites.} La remarque est juste :
la récurrence ci-dessus produit 2.13 en chemin.

\emph{Corollaire 2.14.} $\mathcal{C}$ étant comme ci-dessus, soit
$T : \mathcal{C}^{\circ} \to (\mathrm{Ab})$ un foncteur contravariant exact à
gauche. Soit $(Y_{j})_{j \in I}$ la famille des composantes irréductibles de
$Y = V(I)$, et soit $J \subset I$ une sous-famille telle que, pour tout
$M \in \mathcal{C}$, la nullité de $T(M)$ — qui de toute façon ne dépend que de
$\mathrm{supp}\, M$, en vertu de 2.4 — équivaille à : $Y_{j} \not\subset
\mathrm{supp}\, M$ pour tout $j \in J$. Soit $S$ l'ensemble des points
génériques des $Y_{j}$, $j \in J$. Alors, pour tout $M \in \mathcal{C}$,
\[
  \mathrm{Ass}\, T(M) \;=\; \mathrm{supp}\, M \cap S ;
\]
en particulier le $A$-module $T(M)$ n'a pas d'idéal premier associé
immergé.\footnote{Le feuillet glose $S$ par une addition à l'encre
« $= (Y_{i})_{i \in J}$ », ce qui confond l'ensemble des points génériques avec
la famille des composantes. C'est l'ensemble des points génériques qui est
utilisé partout ensuite, et c'est lui qui est écrit ici.}

Posons $H = \varinjlim_{n} T(A/I^{n})$.\footnote{Le feuillet écrit ici $A/I^{n}$
là où 2.12 écrivait $A_{m} = A/I^{m+1}$. Les deux familles sont cofinales et la
limite est la même ; le décalage d'indice n'a aucune conséquence et n'est pas
corrigé ailleurs.} Par 2.9, $T(M) \simeq \mathrm{Hom}(M,H)$, d'où
\[
  (\ast\ast) \qquad
  \mathrm{Ass}\, T(M) \;=\; \mathrm{supp}\, M \cap \mathrm{Ass}\, H ,
\]
et tout revient à établir $\mathrm{Ass}\, H = S$. On a d'emblée
$\mathrm{Ass}\, H \subset \mathrm{supp}\, H \subset Y$. Soit $y \in Y$ et
$M = A/\mathfrak{p}_{y}$, de sorte que $\mathrm{supp}\, M = \overline{\{y\}}$.
Par $(\ast\ast)$, $T(M) \neq 0$ équivaut à ce que $\overline{\{y\}}$ rencontre
$\mathrm{Ass}\, H$ ; par hypothèse, $T(M) \neq 0$ équivaut à ce que
$\overline{\{y\}}$ rencontre $S$. Donc, pour tout $y \in Y$ :
\[
  \overline{\{y\}} \cap \mathrm{Ass}\, H \neq \emptyset
  \quad\Longleftrightarrow\quad
  \overline{\{y\}} \cap S \neq \emptyset .
\]
Prenons $y \in \mathrm{Ass}\, H$ : il existe alors $z \in S$ avec $z \in
\overline{\{y\}}$, c'est-à-dire $\overline{\{z\}} \subset \overline{\{y\}}$ ;
mais $\overline{\{z\}}$ est une composante irréductible de $Y$ et
$\overline{\{y\}} \subset Y$, donc $\overline{\{y\}} = \overline{\{z\}}$ et
$y = z \in S$. D'où $\mathrm{Ass}\, H \subset S$. Réciproquement, pour $z \in S$,
$\overline{\{z\}}$ rencontre $S$, donc rencontre $\mathrm{Ass}\, H$ en un point
$w$ ; par ce qui précède $w \in S$, et deux éléments de $S$ dont l'un spécialise
l'autre sont égaux, donc $w = z$ et $z \in \mathrm{Ass}\, H$.\footnote{La
dernière phrase du feuillet — « puis $S \subset \mathrm{Ass}\, H$ puisque tout
$y \in S$ doit avoir une spécialisation dans $S$, donc être dans $S$ » — est
elliptique et, prise au pied de la lettre, circulaire. L'argument restitué
ci-dessus est celui qu'elle abrège, et son ressort est bien la dernière clause
qu'il écrit : deux éléments de $S$ dont l'un spécialise l'autre sont
identiques.} L'énoncé dit, en langage d'aujourd'hui, qu'un module de
$I$-torsion représentant un foncteur exact à gauche est une somme d'enveloppes
injectives indexée par des points génériques : pas de premier immergé.

\pagerange{11}{13}
\subsection*{Injectivité, et la généralisation aux faisceaux (pages 11 à 13)}

\emph{Proposition 2.15.} Avec les notations précédentes, pour que $H$ soit un
$A$-module injectif, il faut et il suffit que le foncteur
$T_{H}(M) = \mathrm{Hom}(M,H)$ sur la catégorie $\mathcal{C}$ des $A$-modules de
type fini annulés par une puissance de $I$ soit exact à droite, donc exact.
Autrement dit, 2.10 induit une équivalence entre la catégorie des foncteurs
contravariants \emph{exacts} de $\mathcal{C}$ dans $(\mathrm{Ab})$ et la
catégorie des $A$-modules injectifs dont tout élément est annulé par une
puissance de $I$.

La nécessité est triviale. Pour la suffisance, il suffit — c'est le critère de
Baer, qu'il cite sous le nom de \emph{Tohoku} — de vérifier que pour tout
sous-module $N$ de $A$, l'application $\mathrm{Hom}(A,H) \to \mathrm{Hom}(N,H)$
est surjective, ce qui ramène au cas où $M$ et $N$ sont de type fini. Soit donc
$u : N \to H$. Son image est un sous-module de type fini de $H$, donc annulé par
une puissance de $I$ : $u$ est nul sur $I^{n}N$. Par le lemme d'Artin--Rees, il
existe $m$ tel que $N \cap I^{m}M \subset I^{n}N$.\footnote{Le feuillet écrit
« En vertu de \emph{Krull}, $I^{n}N$ contient une intersection $N \cap I^{m}N$ »
— avec $N$ au second membre, là où le déploiement qui suit immédiatement écrit
$N \cap I^{m}M$, qui est le seul terme dont l'énoncé ait besoin. C'est le lemme
d'Artin--Rees ; le théorème d'intersection de Krull en est le corollaire usuel,
et c'est lui qu'il nomme.} On a alors une injection
$0 \to N/(N \cap I^{m}M) \to M/I^{m}M$ dans $\mathcal{C}$, et $u$ se factorise
par un homomorphisme $u' : N/(N \cap I^{m}M) \to H$, lequel par hypothèse se
prolonge à $M/I^{m}M$ ; d'où $v : M \to H$ prolongeant $u$.

\emph{Corollaire 2.16.} Soit $K$ un $A$-module et soit $H^{0}_{I}(K)$ le
sous-module des éléments annulés par une puissance de $I$. Si $K$ est injectif,
$H^{0}_{I}(K)$ l'est aussi.

En effet, lorsque $M$ parcourt $\mathcal{C}$, l'inclusion induit un isomorphisme
fonctoriel $\mathrm{Hom}(M, H^{0}_{I}(K)) \xrightarrow{\ \sim\ }
\mathrm{Hom}(M,K)$ — un homomorphisme d'un module de $I$-torsion de type fini
atterrit dans la partie de $I$-torsion — et il suffit d'appliquer 2.15.
C'est l'énoncé que le foncteur de $I$-torsion préserve les
injectifs.\footnote{La notation $H^{0}_{I}$ est la sienne, et c'est la seule
occurrence, dans tout le tapuscrit, d'un symbole de cohomologie locale. Le
foncteur est aujourd'hui noté $\Gamma_{I}$ et ses dérivés $H^{i}_{I}$ ; le
module $H$ de 2.9 et les $H^{i}$ de 2.12 sont exactement ces dérivés, et c'est
par là que le tapuscrit rejoint la seconde moitié du dossier. Le feuillet ne le
dit nulle part.}

Juste avant la dernière remarque, une ligne entière est biffée à la machine :
\emph{« Si on rejette tout ce numéro (comme il me semble raisonnable) »}. Le
numéro n'a pas été rejeté ; il est dans la chemise.

\emph{Remarques 2.17.} On peut généraliser 2.6, 2.9 et leurs corollaires au cas
d'un préschéma noethérien $X$, $\mathcal{C}$ désignant la catégorie des modules
cohérents sur $X$ à support contenu dans $Y$ — avec $Y = X$ dans le cas de 2.6.
Pour un foncteur contravariant exact à gauche $T$ de $\mathcal{C}$ dans
$(\mathrm{Ab})$, on a une représentation, essentiellement unique,
\[
  T \;\simeq\; \big(F \rightsquigarrow \mathrm{Hom}(F,H)\big) ,
  \qquad H \simeq \varinjlim_{n} T(\mathcal{O}_{X}/I^{n}) ,
\]
où $H$ est quasi-cohérent à support contenu dans $Y$. La démonstration tient en
deux observations : $T$ est ind-représentable, par le critère général de
\emph{Technique de descente et théorèmes d'existence en géométrie algébrique,
II}, applicable ici parce que les objets de $\mathcal{C}$ sont
noethériens ;\footnote{Le numéro de l'exposé est laissé en blanc sur le
feuillet, quatre points de suspension. Il s'agit du Séminaire Bourbaki,
exposé 195 (1959/60), sous le titre « Le théorème d'existence en théorie
formelle des modules » ; le numéro est fourni ici et ne figure pas sur la page.}
et $\mathrm{Ind}(\mathcal{C})$ est naturellement équivalente à la catégorie de
tous les modules quasi-cohérents sur $X$ à support contenu dans $Y$ — ce qui est
la forme catégorique du fait qu'un module quasi-cohérent est limite inductive
filtrante de modules cohérents. Il conclut : « Nous n'aurons sans doute pas
besoin de ce fait dans ce Chapitre, où les résultats affines développés ici
suffiront pour notre propos. »

\pagerange{15}{20}
\section*{II. Les faisceaux d'ensembles : $(P_1)$ et $(P_2)$ (pages 15 à 20)}

Ici la main change, et le régime d'écriture avec elle. Ce qui suit est un
brouillon rapide : les énoncés et les formules affichées passent, les phrases
qui les relient très souvent non. Le lecteur doit savoir, avant d'entrer, que
les pages 17, 19 et 20 sont en grande partie illisibles et que rien n'y est
comblé ici.\footnote{La chemise de kraft qui enveloppe ces six feuillets porte,
biffés de trois longues diagonales : « Foncteurs élémentaires / sur les espaces
topologiques, / en particulier noethériens, / profondeur $\geqslant 2$ / pour
faisceaux / d'ensembles ». C'est le seul endroit du dossier où le mot
\emph{profondeur} est écrit, et il nomme ce que $(P_{2})$ signifie.}

\emph{Définition 1.} Soient $X$ un espace topologique, $Y$ une partie fermée,
$U = X - Y$, $i : U \to X$ l'inclusion, et $F$ un faisceau d'ensembles sur $X$.
On dit que $F$ a la propriété $(P_{1})$, resp. $(P_{2})$, si le morphisme
canonique
\[
  F \longrightarrow i_{*}i^{*}F \;=\; i_{*}(F|U)
\]
est injectif, resp. bijectif. De façon équivalente, si pour tout ouvert
$V \subset X$ l'application de restriction
\[
  \Gamma(V,F) \longrightarrow \Gamma(V \cap U, F)
\]
est injective, resp. bijective.

En marge, et sans suite : « N.B. regarder aussi la propriété :
\emph{surjective} ». La question restera sans réponse dans cette partie ; elle
est exactement celle à laquelle la troisième partie du dossier répondra, sous le
nom de $(P_{3})$ et pour des objets au lieu de sections.

Ce qui suit — la restriction de l'étude aux espaces localement noethériens dont
les parties fermées se comportent bien, l'introduction de ce qu'il appelle un
« espace noethérien spécial », et les définitions de $\mathrm{Loc}_{x}(X)$ et de
$\mathrm{Loc}_{x}(F)$ — est annoncé sur le feuillet 15 dans des lignes dont
presque rien n'a été lu. La \emph{Définition 2}, qui y transporte $(P_{1})$ et
$(P_{2})$ en un point, est dans le même état, et y nomme au passage des
propriétés $(P_{3})$ et $(P_{4})$ qui n'apparaîtront plus jamais dans cette
partie.\footnote{L'état du feuillet est tel que la transcription y porte, sur
trois lignes consécutives, sept marques d'illisibilité et trois lectures
douteuses. Aucune reconstitution n'est tentée ici : ni la définition de
« l'espace noethérien spécial », ni celle de $\mathrm{Loc}_{x}$, ni le contenu
de $(P_{3})$ et $(P_{4})$ à ce stade. Ce qui suit utilise la seule définition
que la suite du dossier rende certaine, et qui est écrite au net au Lemme 5 :
$X_{(x)}$ est l'ensemble des générisations de $x$ muni de la topologie induite,
et $F_{(x)}$ la restriction de $F$ à cet ensemble.}

La définition ponctuelle qui est effectivement utilisée dans tout le reste du
dossier est donc celle-ci : $F$ a la propriété $(P_{1})$, resp. $(P_{2})$,
\emph{en un point} $x$ lorsque $F_{(x)}$ l'a sur $X_{(x)}$ relativement à
$\{x\}$ ; et sur un espace localement noethérien quelconque, $F$ a la propriété
lorsqu'il l'a en tout point.

\emph{Proposition 3.} Soient $X$ un espace localement noethérien
« spécial »,\footnote{Le terme est le sien et n'est défini nulle part
lisiblement dans le dossier. Ce que les démonstrations en utilisent — toute
partie fermée irréductible a un point générique unique, les parties
constructibles se comportent bien vis-à-vis des générisations, tout fermé non
vide a un point maximal — est ce qu'on appelle aujourd'hui un \emph{espace
spectral}, terme dû à Hochster (1969), soit postérieur au dossier. On ne trouve
sur les feuillets ni ce mot ni celui de \emph{sobre}.} $Y$ une partie fermée,
$F$ un faisceau d'ensembles sur $X$. Sont équivalentes :
\begin{enumerate}
  \item[(i)] $F$ a la propriété $(P_{1})$, resp. $(P_{2})$, relativement à $Y$ ;
  \item[(ii)] pour tout $x \in Y$, $F$ a la propriété $(P_{1})$, resp.
  $(P_{2})$, en $x$.
\end{enumerate}

\emph{Corollaire 4.} Si $F$ satisfait $(P_{1})$, resp. $(P_{2})$, relativement à
$Y$, alors pour tout ouvert $X' \subset X$ et toute partie fermée $Y'$ de $X'$
avec $Y' \subset X' \cap Y$, le faisceau $F|X'$ satisfait $(P_{1})$, resp.
$(P_{2})$, relativement à $Y'$.

La démonstration du corollaire est celle qui passe le mieux sur le feuillet 18,
et elle est purement formelle. Avec $U' = X' - Y'$, on a le triangle commutatif
des restrictions

\begin{tikzcd}
  \Gamma(V,F) \arrow[r, "\rho'"] \arrow[dr, "\rho"'] & \Gamma(V \cap U', F) \arrow[d, "\rho''"] \\
  & \Gamma(V \cap U, F)
\end{tikzcd}

où $\rho$ et $\rho''$ sont injectifs, resp. bijectifs — le second parce que
$V \cap U = (V \cap U') \cap U$ — d'où il suit que $\rho'$ l'est.

\emph{Lemme 5.} Pour tout $x \in X$, en posant $X_{(x)} = \mathrm{Loc}_{x}(X)$,
$X'_{(x)} = X_{(x)} - \{x\}$ et $V_{(x)} = V \cap X_{(x)}$, on a un isomorphisme
\[
  (\ast) \qquad
  \Gamma(V_{(x)}, F_{(x)}) \;\simeq\;
  \varinjlim_{W \ni x} \Gamma(W \cap V, F) ,
\]
la limite étant prise sur les voisinages ouverts $W$ de $x$.\footnote{Le
feuillet 18 écrit la première de ces égalités sous la forme
« $\mathrm{Loc}_{x}(X) = \mathrm{Loc}_{x}(X) - \{x\}$ », le même symbole des
deux côtés : c'est le symbole de gauche qui doit porter un accent, et c'est
$X'_{(x)}$ qui est écrit ici. La transcription note que la formule $(\ast)$ est
la seule ligne de ce feuillet écrite au net, les lignes qui la précèdent étant
un palimpseste de trois états successifs.}

C'est l'énoncé de continuité qui fait fonctionner tout le reste : la
localisation $X_{(x)}$ est l'intersection filtrante décroissante des voisinages
de $x$, et les sections d'un faisceau sur une telle intersection sont la limite
inductive des sections sur les voisinages. C'est de cette forme, et non d'une
autre, que le dossier a besoin, puisque c'est elle qui permet de passer d'une
condition en un point à une condition sur un voisinage.

Le feuillet 17 porte, entre deux ratures, un lemme sans numéro qui est l'outil
topologique du reste : pour deux parties localement constructibles $E \supset
E'$ de $X$ et un point $x$, on a $E \supset E'$ au voisinage de $x$ si et
seulement si
\[
  E_{(x)} \;\supset\; E'_{(x)} ,
\]
ce qui se ramène au fait qu'une partie constructible $Z$ telle que
$Z_{(x)} = \emptyset$ a une adhérence ne contenant pas $x$, c'est-à-dire que
$X - Z$ est un voisinage de $x$. C'est ce lemme qui permet de faire descendre
une égalité vérifiée sur la localisation $X_{(x)}$ jusqu'à un voisinage de
$x$.\footnote{Le feuillet 17 est, avec les feuillets 19 et 20, le plus abîmé du
lot : une longue rature horizontale y traverse quatre lignes, et la prose de
liaison y est presque intégralement illisible. L'énoncé ci-dessus est ce que les
deux formules affichées et les mots « constructible », « immédiat » et
« voisinage » permettent de restituer ; le reste du feuillet, y compris la
démonstration de surjectivité qu'il amorce, n'est pas restitué ici et ne doit
pas l'être.}

Le feuillet 19 porte la démonstration de (ii) $\Rightarrow$ (i), et c'est le
feuillet le plus travaillé du lot : trois longues ratures horizontales, deux
blocs cernés, un mot biffé en marge. Ce qui en passe est la \emph{structure} et
non le détail, et on n'en dira ici que la structure. La question étant locale
sur $X$, on peut supposer $X$ noethérien, et l'on procède par \emph{récurrence
noethérienne} sur les parties fermées : on suppose la propriété acquise
relativement à tout fermé $Y' \subsetneq Y$ et on l'établit relativement à $Y$.
Le cas $Y = \emptyset$ est trivial. Il faut alors montrer que
\[
  \Gamma(X,F) \longrightarrow \Gamma(U,F) , \qquad U = X - Y ,
\]
est injective, resp. bijective. Pour l'injectivité, deux sections $f$, $g$ de
$F$ sur $X$ de même image dans $\Gamma(U,F)$ coïncident en tout point $x$ de $Y$
dont l'adhérence n'est pas $Y$ tout entier, par l'hypothèse de récurrence ; en
un point maximal de $Y$, c'est l'hypothèse $(P_{1})$ en ce point qui conclut, via
le Lemme 5. La surjectivité, au feuillet 20, suit le même plan : on recouvre
$V_{(x)}$ par des ouverts $W_{i}$ portant des sections $f_{i}$ qui induisent la
section donnée, on considère les parties
$W_{ij} = \{\, y \in W_{i} \cap W_{j} : f_{i,y} = f_{j,y} \,\}$, qui sont
constructibles, et l'on ajuste les $W_{i}$ pour que ces parties soient les
intersections entières, de sorte que les $f_{i}$ se recollent en une vraie
section.\footnote{Le feuillet 20 écrit « $f_{i,y} = h_{y}$ » là où l'argument
demande $f_{j,y}$ : la transcription porte le $h$ tel quel et le signale. La
lecture retenue ici est $f_{j}$, seule compatible avec le rôle de ces parties
dans le recollement. Il faut redire que ces deux feuillets sont, en prose, très
largement illisibles : le décompte des étapes ci-dessus est ce que les formules
affichées et les mots isolés qui subsistent permettent d'affirmer, et rien de
plus n'est affirmé.}

L'implication (i) $\Rightarrow$ (ii) occupe la fin du feuillet 20 et est plus
courte : pour $x \in Y$, le Corollaire 4 appliqué à $\overline{\{x\}}$ donne que
$F$ satisfait la propriété relativement à $\overline{\{x\}}$, donc, avec
$V = X - \overline{\{x\}}$, que $\Gamma(W,F) \to \Gamma(W \cap V, F)$ est
injective, resp. bijective, pour tout voisinage ouvert $W$ de $x$ ; on passe à
la limite sur $W$ et l'on identifie le résultat, par deux applications du
Lemme 5, à
\[
  \Gamma(X_{(x)}, F_{(x)}) \longrightarrow \Gamma(X'_{(x)}, F_{(x)}) .
\]

\pagerange{21}{28}
\section*{III. Les catégories fibrées : $(P_1)$, $(P_2)$, $(P_3)$ (pages 21 à 28)}

\pagerange{21}{24}
\subsection*{Les trois conditions, et leur forme ponctuelle (pages 21 à 24)}

Le feuillet 21 s'ouvre sur la fin de la démonstration commencée au feuillet 20 —
le prolongement d'une section de $F$ sur $U$ en une section sur $X$, par
l'hypothèse de récurrence et la propriété $(P_{2})$ au point $x$ — puis passe
sans transition au changement de cadre qui occupe les huit derniers feuillets.

Le changement est celui-ci. Jusqu'ici on prolongeait des \emph{sections}, et
deux sections sont égales ou ne le sont pas. On va prolonger des \emph{objets},
et deux objets ne sont pas égaux : ils sont isomorphes, et en général de
plusieurs manières. La donnée pertinente n'est plus un faisceau mais une
\emph{catégorie fibrée} $\mathcal{F}$ sur la catégorie des ouverts de $X$ : une
catégorie $\mathcal{F}(V)$ pour chaque ouvert, des foncteurs de restriction, et
les isomorphismes de compatibilité entre eux. Et un foncteur de restriction peut
manquer d'être une équivalence de trois manières emboîtées, d'où trois
conditions au lieu de deux.

\emph{Définition 6.} Soient $X$ un espace topologique, $Y$ une partie fermée,
$U = X - Y$, et $\mathcal{F}$ une catégorie fibrée sur la catégorie des ouverts
de $X$. On dit que $\mathcal{F}$ satisfait la condition $(P_{1})$, resp.
$(P_{2})$, resp. $(P_{3})$, relativement à $Y$, si pour tout ouvert
$V \subset X$ le foncteur de restriction
\[
  \mathcal{F}(V) \longrightarrow \mathcal{F}(V \cap U)
\]
est fidèle, resp. pleinement fidèle, resp. une équivalence de
catégories.\footnote{Le feuillet porte devant le $\mathcal{F}$ du premier membre
un jambage supplémentaire, lu $\Gamma$ par la transcription, que le second
membre n'a pas ; l'asymétrie est sur la page et la lecture y est donnée pour
douteuse. Aucune interprétation ne rend $\Gamma\mathcal{F}(V)$ différent de
$\mathcal{F}(V)$ dans cet énoncé, et c'est $\mathcal{F}(V)$ qui est écrit ici.}
On dit que $\mathcal{F}$ satisfait la condition en un point $x \in X$ si le
foncteur
\[
  \mathrm{Pseudolim}_{W}\ \mathcal{F}(W)
  \longrightarrow
  \mathrm{Pseudolim}_{W}\ \mathcal{F}(W - W \cap \overline{\{x\}})
\]
est fidèle, resp. pleinement fidèle, resp. une équivalence, $W$ parcourant les
voisinages ouverts de $x$.\footnote{Le mot est le sien, écrit avec le $W$ sous
une longue flèche : c'est une limite inductive filtrante de \emph{catégories},
prise au sens où la composition des restrictions n'est associative qu'à
isomorphisme près. On dit aujourd'hui pseudo-colimite filtrante, ou
$2$-colimite. Le nom de « limite » est trompeur : $W$ décroît, la construction
est une colimite.}

Une note marginale, elle-même très abîmée, demande si les deux formes — globale
et ponctuelle — coïncident sur un espace localement noethérien spectral. C'est
la Proposition 7, et c'est la question que toute cette partie traite.

Les pages 22 et 23 ramènent les trois conditions à des conditions sur des
faisceaux, ce qui est ce qui permettra à la deuxième partie du dossier de servir.
Pour deux objets $f$, $g$ de $\mathcal{F}(V)$, soit $\mathcal{H}om(f,g)$ le
faisceau sur $V$ des homomorphismes, dont les sections sur $V' \subset V$ sont
les homomorphismes $f|V' \to g|V'$. Alors :
\begin{itemize}
  \item $\mathcal{F}$ satisfait $(P_{1})$, resp. $(P_{2})$, relativement à $Y$
  si et seulement si, pour tout ouvert $V$ et tout couple
  $f, g \in \mathrm{Ob}\,\mathcal{F}(V)$, le faisceau $\mathcal{H}om(f,g)$
  satisfait $(P_{1})$, resp. $(P_{2})$, relativement à $Y \cap V$ ;
  \item $\mathcal{F}$ satisfait $(P_{3})$ si de plus, pour tout ouvert $V$, tout
  objet $f'$ de $\mathcal{F}(V \cap U)$ est isomorphe à la restriction d'un objet
  de $\mathcal{F}(V)$.
\end{itemize}

\footnote{Le feuillet écrit ce dernier objet dans $\mathcal{F}(W)$ alors que
l'argument demande $\mathcal{F}(V)$ ; la transcription le signale et porte $W$
tel quel. C'est $V$ qui est écrit ici. Une hypothèse est par ailleurs
tacitement faite dans toute cette page, et il faut la dire : pour que
« $V' \rightsquigarrow \mathrm{Hom}(f|V', g|V')$ » soit un \emph{faisceau} et
non un simple préfaisceau, il faut que $\mathcal{F}$ soit un préchamp. La
Définition 6 ne demande qu'une catégorie fibrée.}

La troisième forme, à la page 23, est celle qui rend la condition ponctuelle
maniable : pour tout $x \in X$, on introduit la \emph{catégorie fibrée induite}
$\mathcal{F}_{(x)}$ sur $X_{(x)}$ en posant
\[
  \mathcal{F}_{(x)}(w) \;=\; \mathrm{Pseudolim}_{W}\ \mathcal{F}(W) ,
\]
$W$ parcourant les ouverts de $X$ tels que $W_{(x)} = w$. On a alors, pour $f$,
$g$ au-dessus d'un ouvert $V$,
\[
  \mathcal{H}om(f,g)_{(x)} \;\simeq\; \mathcal{H}om(f_{(x)}, g_{(x)}) ,
\]
et le N.B. du feuillet en tire ce qui sera utilisé partout : la propriété
$(P_{1})$, resp. $(P_{2})$, resp. $(P_{3})$, de $\mathcal{F}$ \emph{en} $x$
revient à la même propriété de $\mathcal{F}_{(x)}$ \emph{sur} $X_{(x)}$, et pour
$(P_{1})$ et $(P_{2})$ elle revient encore à ce que les faisceaux
$\mathcal{H}om(f,g)$ aient la propriété correspondante en $x$.\footnote{Le
feuillet écrit cette identification avec une double flèche et non une flèche
d'isomorphisme, ce que la transcription porte tel quel ; il l'introduit aussi
par « j'imagine », qui est sa manière d'annoncer ce qu'il n'a pas vérifié. La
lecture la donne comme isomorphisme, ce qu'elle est, la limite inductive
filtrante commutant aux Hom d'objets de présentation finie sur un voisinage
fixé.}

\emph{Proposition 7.} Soient $X$ un espace topologique noethérien spectral,
$\mathcal{F}$ une catégorie fibrée sur la catégorie des ouverts de $X$, et $Y$
une partie fermée. Sont équivalentes :
\begin{enumerate}
  \item[(i)] $\mathcal{F}$ satisfait $(P_{1})$, resp. $(P_{2})$, resp.
  $(P_{3})$, relativement à $Y$ ;
  \item[(ii)] pour tout $x \in Y$, $\mathcal{F}$ satisfait $(P_{1})$, resp.
  $(P_{2})$, resp. $(P_{3})$, en $x$.
\end{enumerate}

\footnote{La condition (ii) est écrite sur le feuillet « en $X$ », avec une
majuscule, là où c'est le point $x$ qui est en cause ; la transcription le
signale. Pour $(P_{1})$ et $(P_{2})$, l'énoncé n'est autre que la Proposition 3
appliquée aux faisceaux $\mathcal{H}om(f,g)$, ce que la page dit en une ligne :
« Reste le cas $P_{3}$, quand $P_{1}$, $P_{2}$ sont déjà vérifiés. »}

\emph{Corollaire 8.} Si $\mathcal{F}$ satisfait $(P_{1})$, resp. $(P_{2})$,
resp. $(P_{3})$, relativement à $Y$, alors pour tout ouvert $X' \subset X$ et
toute partie fermée $Y' \subset Y \cap X'$, la catégorie fibrée induite
$\mathcal{F}|X'$ satisfait la même condition relativement à $Y'$.

\pagerange{25}{27}
\subsection*{Les démonstrations, et la traduction cohomologique (pages 25 à 27)}

Le Corollaire 8 se démontre comme le Corollaire 4, et par le même diagramme, une
catégorie en remplaçant un ensemble de sections :

\begin{tikzcd}
  \mathcal{F}(V) \arrow[r, "\rho'"] \arrow[dr, "\rho"'] & \mathcal{F}(V \cap U') \arrow[d, "\rho''"] \\
  & \mathcal{F}(V \cap U)
\end{tikzcd}

$\rho$ et $\rho''$ étant des équivalences, $\rho'$ en est une.

L'implication (i) $\Rightarrow$ (ii) de la Proposition 7 s'obtient en appliquant
la condition à $\overline{\{x\}}$ pour $x \in Y$ : pour tout voisinage ouvert
$W$ de $\overline{\{x\}}$, le foncteur $\mathcal{F}(W) \to \mathcal{F}(W - W
\cap \overline{\{x\}})$ est une équivalence, donc le reste après passage à la
limite sur les $W$, ce qui est exactement la condition en $x$.\footnote{Le
feuillet abrège les trois énoncés — fidélité, pleine fidélité, équivalence — par
une lettre $\mathcal{H}$ qu'il n'introduit nulle part et qu'il glose seulement
par « l'un des types ». Elle n'est pas reprise ici ; les trois cas sont énoncés
ensemble, comme partout ailleurs dans le dossier.}

L'implication (ii) $\Rightarrow$ (i) est le seul endroit qui coûte quelque
chose, et seulement pour $(P_{3})$ : il s'agit de montrer que tout objet $f'$ de
$\mathcal{F}(U)$ provient d'un objet de $\mathcal{F}(X)$. La question étant
locale, on suppose $X$ noethérien, et l'on fait la récurrence noethérienne :
$\mathcal{F}$ satisfait $(P_{3})$ relativement à tout fermé $Y' \subsetneq Y$.
Si $Y = \emptyset$, il n'y a rien à prouver. Sinon, soit $x$ un point maximal de
$Y$. La restriction $f'|X'_{(x)}$ provient, par hypothèse, d'un objet de
$\mathcal{F}_{(x)}(X'_{(x)})$, donc d'un objet $f_{0}$ de $\mathcal{F}(W)$ pour
un voisinage ouvert $W$ de $x$ assez petit. Les restrictions de $f_{0}$ et de
$f'$ à $W \cap U$ deviennent isomorphes sur $V_{(x)} = X'_{(x)}$, donc — quitte
à rétrécir $W$ — sont isomorphes. En recollant le long de cet isomorphisme, on
obtient un objet $f'_{1}$ de $\mathcal{F}(W \cup U)$ qui prolonge $f'$ ; le
complémentaire de $W \cup U$ est un fermé strictement contenu dans $Y$, donc
l'hypothèse de récurrence fait provenir $f'_{1}$ d'un objet de
$\mathcal{F}(X)$.\footnote{C'est le seul pas du dossier qui demande à
$\mathcal{F}$ plus que d'être une catégorie fibrée. Recoller $f_{0}$ sur $W$ et
$f'$ sur $U$ le long d'un isomorphisme au-dessus de $W \cap U$, c'est
l'effectivité de la descente pour le recouvrement ouvert $\{W, U\}$ :
$\mathcal{F}$ doit être un \emph{champ}. La Définition 6 ne le demande pas, et
aucun feuillet ne le dit. L'hypothèse est ajoutée ici parce que sans elle la
démonstration ne tient pas ; le vocabulaire — champ, préchamp, descente
effective — est postérieur au dossier.}

\emph{Corollaire 10.} Sont équivalentes :
\begin{enumerate}
  \item[(i)] les conditions $(P_{1})$, resp. $(P_{2})$, resp. $(P_{3})$ ;
  \item[(ii)] $\underline{H}^{i}_{Y}(F) = 0$ pour $i \leqslant 0$, resp.
  $i \leqslant 1$, resp. $i \leqslant 2$.
\end{enumerate}

\footnote{C'est ici que la lecture s'écarte le plus du feuillet, et il faut dire
exactement en quoi. Le feuillet écrit les trois plages « $i \leqslant 1$
($i \leqslant 2$, $i \leqslant 3$) », soit une de trop dans chacun des trois
cas. Deux choses le montrent, et elles sont toutes deux sur les feuillets. La
première est la Proposition 9 du feuillet suivant, qui énonce les deux premiers
cas sous la forme « $\underline{H}^{0}_{Y}(F) = 0$
($\underline{H}^{0}_{Y}(F) = 0$ et $\underline{H}^{1}_{Y}(F) = 0$), i.e. $i = 0$
($i = 0, 1$) » — ce qui est $i \leqslant 0$ et $i \leqslant 1$, non
$i \leqslant 1$ et $i \leqslant 2$. La seconde est le déploiement que le
feuillet 27 donne lui-même immédiatement après, pour traiter « le cas
$P_{3}$ » : il y écrit $\underline{H}^{2}_{Y}$, et non $\underline{H}^{3}_{Y}$.
Le feuillet se corrige donc deux fois lui-même, et ce sont ses corrections qui
sont retenues. Une remarque sur l'objet de l'énoncé : (i) y est écrit « $F$
satisfait à $P_{1}$ ($P_{2}$, $P_{3}$) », avec $F$ là où toute la page précédente
porte sur $\mathcal{F}$. Les trois conditions n'ont de sens à trois niveaux que
pour une catégorie fibrée ; c'est ainsi que l'énoncé est lu ici, $F$ étant le
faisceau abélien dont $\mathcal{F}$ est, comme au feuillet suivant, la catégorie
des torseurs.}

Le feuillet donne, pour le cas $(P_{3})$, l'identification qui fait tout :
\[
  \underline{H}^{2}_{Y}(F) \;=\;
  \text{faisceau associé au préfaisceau } \;
  V \rightsquigarrow H^{1}(V \cap U, F) .
\]
C'est le cas $i = 2$ de la formule générale $\underline{H}^{i}_{Y}(F) \simeq
R^{i-1}i_{*}(F|U)$ pour $i \geqslant 2$, et c'est elle qui donne l'implication
(i) $\Rightarrow$ (ii). Réciproquement, si les faisceaux s'annulent, tout objet
$f'$ de $\mathcal{F}(U)$ provient d'un objet de $\mathcal{F}(X)$, la question
étant locale sur $Y$ et réglée en chaque point par la Proposition
7.\footnote{L'exposant du $H$ du second membre ne se lit pas sur le feuillet :
la transcription note « un seul jambage, ni un zéro net ni un chiffre », et
celui de $\underline{H}_{Y}$ s'y lit $2$. Le $1$ écrit ici est notre lecture, et
il est forcé : avec $\underline{H}^{2}_{Y}$ au premier membre, la seule valeur
qui rende la formule vraie est $1$, puisque $\underline{H}^{i}_{Y}(F)$ est le
faisceau associé à $V \rightsquigarrow H^{i-1}(V \cap U, F)$ pour
$i \geqslant 2$. Un seul jambage est d'ailleurs ce à quoi ressemble un 1. Les
deux autres occurrences du même exposant, dans les deux lignes qui suivent, sont
illisibles de la même façon et sont lues de la même manière.}

\pagerange{28}{28}
\subsection*{Le cas des torseurs (page 28)}

Le dernier feuillet donne le cas particulier où tout se calcule, et où le lien
avec la première moitié du dossier devient visible d'un coup.

Soient $F$ un faisceau de groupes sur $X$ et $\mathcal{F}$ la catégorie fibrée
des fibrés principaux homogènes sous $F$ — des \emph{torseurs} — au-dessus des
ouverts de $X$. Soit $Y$ une partie fermée.\footnote{La ligne qui définit
$\mathcal{F}$ est très abrégée sur le feuillet ; la transcription ne donne pour
sûrs que la locution « fibrés principaux homogènes » et le rôle de $F$, et
signale que $F$ y est appelé \emph{préfaisceau}. C'est comme faisceau qu'il est
lu ici, faute de quoi ni $i_{*}i^{*}F$ ni les $\underline{H}^{i}_{Y}$ n'auraient
leur sens usuel.}

\emph{Proposition 9.} Sont équivalentes :
\begin{enumerate}
  \item[(i)] $\mathcal{F}$ satisfait $(P_{1})$, resp. $(P_{2})$, relativement à
  $Y$ ;
  \item[(ii)] $F$ satisfait $(P_{1})$, resp. $(P_{2})$, relativement à $Y$ ;
  \item[(iii)] $\underline{H}^{0}_{Y}(F) = 0$, resp.
  $\underline{H}^{0}_{Y}(F) = 0$ et $\underline{H}^{1}_{Y}(F) = 0$.
\end{enumerate}

La démonstration tient en trois phrases et elles sont toutes les trois sur le
feuillet. (ii) $\Rightarrow$ (i) : si $f_{0}$ désigne le torseur trivial sur
$X$, on a $F \simeq \mathcal{H}om(f_{0}, f_{0})$, et la traduction du feuillet 22
s'applique. (i) $\Rightarrow$ (ii) : pour $f$, $g$ dans $\mathcal{F}(V)$, le
faisceau $\mathcal{H}om(f,g)$ est localement isomorphe à $F$, donc hérite de la
propriété. (ii) $\Leftrightarrow$ (iii) : par définition même de
$\underline{H}^{0}_{Y}(F)$ et de $\underline{H}^{1}_{Y}(F)$ comme noyau et
conoyau de
\[
  F \longrightarrow i_{*}i^{*}F ,
\]
l'injectivité de ce morphisme est l'annulation du premier et sa bijectivité
l'annulation des deux.\footnote{La suite exacte à quatre termes
$0 \to \underline{H}^{0}_{Y}(F) \to F \to i_{*}i^{*}F \to
\underline{H}^{1}_{Y}(F) \to 0$ est ce qui justifie l'appellation, et elle rend
l'équivalence immédiate. Deux réserves, qu'aucun feuillet ne lève. La première :
le conoyau n'est un faisceau de groupes que si $F$ est abélien ; pour $F$ non
abélien, ce que l'énoncé appelle $\underline{H}^{1}_{Y}(F)$ est un faisceau
d'ensembles pointés, quotient de $i_{*}i^{*}F$ par $F$, et l'énoncé garde son
sens à cette lecture près. La seconde : il n'y a pas ici de troisième condition.
$(P_{3})$ pour la catégorie des torseurs demanderait qu'un torseur sur $U$ se
prolonge à $X$, ce qui relève de $\underline{H}^{2}_{Y}$ par le feuillet
précédent, et la Proposition 9 ne traite que $(P_{1})$ et $(P_{2})$. Le dossier
s'arrête là.}

Et c'est ici qu'il faut revenir au tapuscrit. Les faisceaux
$\underline{H}^{i}_{Y}(F)$ dont ce dernier feuillet fait le critère sont, sur le
spectre d'un anneau noethérien, exactement les faisceaux associés aux modules
\[
  H^{i} \;=\; \varinjlim_{m} \mathrm{Ext}^{i}_{A}(A/I^{m+1}, F)
\]
que le Corollaire 2.12 construit à la page 10, comme limite inductive des
valeurs d'un foncteur cohomologique contravariant sur les $A_{m}$. Le tapuscrit
démontre que cette limite représente le foncteur, qu'elle suffit à tester
l'annulation, et que ses idéaux premiers associés sont les points génériques des
composantes de $Y$ ; le manuscrit s'en sert pour dire quand un objet défini hors
de $Y$ se prolonge. Aucun des vingt-quatre feuillets ne dit que ce sont les
mêmes $H^{i}$. C'est pourtant la seule raison pour laquelle ils sont dans la
même chemise.

\end{document}
